\subsection{选题侧重与成员分工}

从近六年A、B、C题的任务差异看，A题常要把物理过程写成方程并处理数值误差，
B题更看重决策变量、约束和可行方案，C题则从数据口径、预测或评价进入决策。
据此，队长负责机理与连续模型，组员1负责规划和优化，组员2负责数据与统计学习。
三人共用同一套输入输出表：队长给出状态量和物理边界，组员1接收状态量建立目标
与约束，组员2整理参数、特征和误差。这样分工后，一项结果由谁计算、交给谁使用，
在代码和论文中都能找到对应位置。

本阶段按三项标准记录学习成果：能用自己的话说清模型为何适用；能把公式写成程序
步骤；能指出它在什么条件下不该使用。思维导图中的方法很多，本次先选十三个单元，
以后再用真题补足尚未动手的分支。

